% O014 / D80 — historical pre-introduction notes, English edition.
% Complete translation of pre-prelude.tex © Wen-Wei Li, CC BY 4.0.
% Stable ID: o014.aljabr2.frontmatter.historical-draft-notes
% This historical source file was not active in the frozen 2024 master build.
\chapter*{Introduction}

\section*{Note on the final draft}
This edition continues the correction of errors and the making of minor
adjustments. The complete version will be released after the publication
process has concluded; please look forward to it. Criticism and corrections
from readers remain most welcome, and necessary changes will be incorporated
into the errata.

\vspace{1em}
\begin{flushright}\begin{minipage}{0.3 \textwidth}
  \begin{tabular}{c}
    Wen-Wei Li \\
    December 2022 \\
    Shijingshan
  \end{tabular}
\end{minipage}\end{flushright}
\vspace{1em}

\section*{Note on the fourth draft}
This draft completes the introductions at the beginning of each chapter and
also concludes a round of checking of the main text and appendices. The result
was the removal of numerous typographical and mathematical errors, the
improvement of awkward formulations, and several structural adjustments. The
next circulated version is expected to be the finished version, including the
introduction, although the present content may still be revised.

\vspace{1em}
\begin{flushright}\begin{minipage}{0.3 \textwidth}
  \begin{tabular}{c}
    Wen-Wei Li \\
    November 2022 \\
    northern Haidian
  \end{tabular}
\end{minipage}\end{flushright}
\vspace{1em}

\section*{Note on the third draft}
Compared with the preceding draft, this edition inserts Chapter Six on group
homology and cohomology, adds appendices, and moves some earlier main-text
material into the appendices. Existing material has also been cleaned up,
corrected, and improved. This edition adds to the appendices a somewhat
circuitous proof of the Freyd--Mitchell embedding theorem using the technique
of $\IndC$-completion, even though the main text makes almost no use of the
embedding theorem.

The overall framework of the book has essentially taken shape, but the
introduction and the opening sections of every chapter still need to be
written. As the note on the second draft observed, this is a major omission.

Moreover, the entire book still requires a complete review. This edition
undoubtedly contains many errors and many passages that remain to be put in
order.

\vspace{1em}
\begin{flushright}\begin{minipage}{0.3 \textwidth}
  \begin{tabular}{c}
    Wen-Wei Li \\
    August 2022 \\
    northern Haidian
  \end{tabular}
\end{minipage}\end{flushright}
\vspace{1em}

\section*{Note on the second draft}
The present edition consists of eight chapters. In addition to three new
chapters, minor corrections have been made to the first five. According to the
writing plan, other material is quite likely to be inserted or appended in the
future, and existing material may still be adjusted; in particular, some
digressions may need pruning.

It must also be stressed that a mathematical work is by no means merely a heap
of definitions and arguments. Until the introduction and the reading guides at
the beginning of every chapter have actually been completed, every chapter in
this manuscript remains incomplete. Presenting such half-finished work is only
a provisional expedient.

The manuscript inevitably still contains many errors and deficiencies.
Criticism and corrections from readers are earnestly invited.

\vspace{1em}
\begin{flushright}\begin{minipage}{0.3 \textwidth}
  \begin{tabular}{c}
    Wen-Wei Li \\
    March 2022 \\
    Gaobeidian
  \end{tabular}
\end{minipage}\end{flushright}
\vspace{1em}

\section*{Note on the first draft}

The manuscript before the reader is part of the planned Volume Two of
\emph{Methods of Algebra} and was expected to comprise roughly half of the
whole book. Besides the missing later chapters, the introduction and the
reading guidance at the beginning of each chapter are either extremely brief
or absent altogether. The exercise sections are also quite sparse. This is,
without qualification, unfinished work. Many typographical and mathematical
errors are to be expected, and the existing chapters may change in the future.

The main reason for publishing such a draft is my belief that it is useful to
readers and that its audience ought to be broad. The second reason is that the
next milestone remains some time away. Finally, I hope that readers' feedback
will help the finished work appear sooner and in better form.

It is clear that Volume One of \emph{Methods of Algebra} covers all the
algebraic knowledge required here. Since both Volume One and this book use
standard terminology, other textbooks of the usual sort should also provide
the necessary background. Readers are advised first to consult the opening
conventions in order to confirm the notation and conventions in use.

The spirit in which this volume is written is broadly similar to that of
Volume One, although the different subject matter also brings differences in
style. Those differences can be explained carefully only in a future
introduction. Like Volume One, this book does not encourage beginners to read
it sequentially unless they are already sufficiently comfortable with
abstract methods. Concrete reading guidance will appear in the introduction
and at the beginning of each chapter; unfortunately, such guidance can be
written only after the overall form of the book has stabilized.

Many relevant works were consulted extensively during preparation. Some are
already cited in the main text; the remainder were intended to be listed in
the introduction once the whole book was complete. Their omission is not
deliberate, and I ask the reader's understanding.

The present contents correspond roughly to what is traditionally called
homological algebra, which is a proper subset of ``linear algebra'' in its
genuine sense. From a modern point of view, the first five chapters form a
barely adequate system; even as foundational training in classical homological
algebra, however, much remains untouched. The remaining material may have to
await a future Chapter Six. Standard applications such as group cohomology
would have to be postponed to still later chapters.

For the terms governing use of this file, please consult the link on the cover
page. All suggestions, criticism, and corrections are welcome.

\vspace{1em}
\begin{flushright}\begin{minipage}{0.3 \textwidth}
  \begin{tabular}{c}
    Wen-Wei Li \\
    March 2021 \\
    Peking University Library
  \end{tabular}
\end{minipage}\end{flushright}
